% Bản nháp chương 4 theo dàn bài mới, viết với bộ mẫu ở style/mau/.
% Chưa nối vào report.tex.
% ===========================================================================
\chapter{Thuật toán nào về đích trước}
\label{sec:compare}

Chương~\ref{sec:enough} chốt đích cần tới: sai số $10^{-6}$, ứng với sai lệch
giá dưới một đơn vị tiền trên một chiếc máy giá trung vị. Chương này đặt bảy cấu
hình tốt nhất của bảy nhóm thí nghiệm lên cùng một đường chạy tới đích ấy, rồi
đặt cả nhóm cạnh các cài đặt mặc định của scikit-learn.

Câu trả lời gọn nằm ngay ở đây. Newton về đích sau \SI{0.133}{\second}, nhanh hơn
gradient descent 3,5 lần; năm cấu hình tất định còn lại đều cho cùng một RMSE, và
khoảng cách giữa chúng chỉ là chuyện chờ lâu hay chóng. Phần còn lại của chương
nói rõ khoảng cách ấy thay đổi thế nào khi đổi mức chính xác yêu cầu, và vì sao
một cấu hình về đích sớm chưa chắc đi xa được.

\section{Ba cách nhìn cùng một dữ liệu}
\label{sec:compare-figures}

\resultfig{all_methods_iter}
  {So sánh các thuật toán ở cấu hình tốt nhất, theo số vòng lặp. Trục hoành cắt
   ở 200 vòng lặp cho dễ đọc.}
  {fig:all-iter}
\resultfig{all_methods_time}
  {So sánh các thuật toán ở cấu hình tốt nhất, theo thời gian chạy. Trục hoành
   cắt ở 6 giây.}
  {fig:all-time}
\resultfig{rmse_vs_time}
  {RMSE trên tập kiểm tra theo thời gian chạy. Đây là hình trả lời trực tiếp câu
   hỏi của người dùng mô hình: với một ngân sách tính toán cho trước thì phương
   pháp nào cho sai số giá thấp nhất.}
  {fig:rmse-vs-time}

Ba hình cùng dựng từ một bộ dữ liệu nhưng trả lời ba câu khác nhau.
Hình~\ref{fig:all-iter} cho biết mỗi vòng lặp cắt được bao nhiêu sai số, tức nó
đo chất lượng của hướng đi mà thuật toán chọn. Hình~\ref{fig:all-time} cho biết
mỗi vòng lặp ấy đắt bao nhiêu, nên thứ hạng trên nó khác thứ hạng trên hình
trước bất cứ khi nào chi phí mỗi vòng chênh nhau, đúng trường hợp của Newton so
với các phương pháp bậc một. Hình~\ref{fig:rmse-vs-time} bỏ hẳn trục sai số tối
ưu hóa và vẽ thẳng sai số giá theo giây, nên nó đọc được mà không cần biết
$\fstar$ là gì.

\section{Bảng tổng hợp}
\label{sec:compare-table}

\begin{table}[tp]
  \centering
  \caption{Bảy cấu hình tốt nhất, đo trên cùng một máy, mỗi cấu hình chạy ba lần
    lấy trung vị. Cột \emph{tổng thời gian} là thời gian chạy tới khi thuật toán
    dừng, vì hội tụ, vì đình trệ, hoặc vì hết ngân sách vòng lặp.}
  \label{tab:summary}
  \begin{tabular}{llrrr}
    \toprule
    Thuật toán & Cấu hình & Vòng lặp đạt $10^{-6}$ & Thời gian (s) & Tổng (s) \\
    \midrule
    Newton   & $t = 1$, phân rã một lần          & 1         & \num{0.133} & \num{0.133} \\
    L-BFGS   & $m = 20$                          & 8         & \num{0.282} & \num{1.49}  \\
    GD       & $t = 2/L$                         & 20        & \num{0.465} & \num{4.99}  \\
    AGD      & $\beta_k$ tăng dần, có khởi động lại & 20     & \num{0.548} & \num{2.97}  \\
    GD       & backtracking $\alpha = \beta = 0.5$ & 15      & \num{0.563} & \num{7.23}  \\
    SGD      & $B = 64$, giảm bậc thang          & không đạt & --          & \num{2.61}  \\
    SGD      & $B = 256$, bước hằng              & không đạt & --          & \num{1.86}  \\
    \bottomrule
  \end{tabular}
\end{table}

\paragraph{Newton.} Toàn bộ lời giải mất \SI{0.133}{\second}
(bảng~\ref{tab:summary}), tức nhanh hơn 3,5 lần so với thời gian gradient descent
cần chỉ để chạm $10^{-6}$, và nhanh hơn 37 lần so với thời gian nó chạy tới lúc
dừng hẳn. Khoảng cách ấy đến từ chỗ hàm mục tiêu là hàm bậc hai nên bước Newton
chính là nghiệm
đóng \eqref{eq:closed-form}, chứ không đến từ một tính chất nào của phương pháp
áp dụng được cho bài toán khác. Với hàm phi tuyến, Newton cần nhiều vòng lặp và
mỗi vòng vẫn tốn $\bigO(nd^{2} + d^{3})$, nên kết luận này không chuyển sang được.

\paragraph{Các phương pháp bậc một ở mức chính xác vừa phải.} Ba cấu hình GD, AGD
và GD với backtracking chạm $10^{-6}$ trong khoảng \SI{0.465}{\second} đến
\SI{0.563}{\second}, chênh nhau 21\%, và ba đường của chúng bám sát nhau suốt
phần đầu hình~\ref{fig:all-iter}. Giai đoạn đầu do các hướng có độ cong lớn chi
phối, mà mọi phương pháp bậc một đều cắt các hướng ấy nhanh như nhau, nên ở mức
chính xác vừa phải thì chọn phương pháp nào cũng gần như không khác nhau. Chênh
lệch 21\% ấy không nên đọc thành thứ hạng, vì ba lần chạy của cùng một cấu hình
đã dao động 5\% tới 25\% ở phần lớn các nhóm và tới 61\% ở trường hợp xấu nhất.
Chính khoảng dao động đó, chứ không phải bản thân thuật toán, quyết định cấu hình
nào được chọn khi hai cấu hình hòa nhau về số vòng lặp; $t = 2/L$ được chọn thay
cho $t = 2/(L+\mu)$ ở bảng~\ref{tab:summary} đúng theo cách ấy.

\paragraph{Phần đuôi, và chuyện cấu hình tốt nhất phụ thuộc dung sai.} Cấu hình
gradient descent về đích sớm nhất lại là cấu hình không đi xa được. Với
$t = 2/L$, hệ số co rút của thành phần ứng với trị riêng $L$ bằng đúng 1 nên
thành phần đó đứng yên mãi, và lần chạy dừng ở sai số $4{,}9 \cdot 10^{-8}$ thay
vì đi tiếp tới giới hạn số học. Muốn đi tới giới hạn ấy thì phải đổi cấu hình:
$t = 1{,}9/L$ hội tụ hẳn sau \SI{6.04}{\second}, còn $t = 2/(L+\mu)$ chạy hết
ngân sách 1200 vòng trong \SI{40.1}{\second} và dừng ở $8{,}3 \cdot 10^{-16}$. So
với hai con số đó, AGD chạm giới hạn số học sau \SI{2.97}{\second}, tức nhanh hơn
2 lần và 13,5 lần. Lợi thế của gia tốc vì thế chỉ hiện ra ở phần đuôi, nơi $\mu$
chi phối tốc độ hội tụ chứ không còn là các hướng có độ cong lớn.

Nói cách khác, câu "cấu hình nào tốt nhất" không có nghĩa nếu chưa nói tốt cho
mức chính xác nào. Ở dung sai của chương~\ref{sec:enough} thì $t = 2/L$ là lựa
chọn đúng, và việc nó không xuống được dưới $5 \cdot 10^{-8}$ hoàn toàn không
quan trọng vì mức ấy đã nằm sâu dưới ngưỡng mà bài toán định giá phân biệt được.

\paragraph{Hai cấu hình SGD.} Cả hai đều không đạt $10^{-6}$ trong ngân sách 30
lượt duyệt dữ liệu, nhưng chúng dừng ở hai chỗ rất khác nhau. Bản giảm theo bậc
thang dừng ở $5{,}8 \cdot 10^{-6}$, tức chỉ lệch \num{0.6} đơn vị tiền so với
nghiệm đóng, còn bản bước hằng dừng ở $8{,}2 \cdot 10^{-4}$ và lệch \num{87.7}
đơn vị. Chênh lệch giữa hai bản nằm ở quy tắc chọn độ dài bước chứ không ở kích
thước lô, và chương~\ref{sec:mechanisms} truy nguyên nhân.

\section{Đặt cạnh thư viện}
\label{sec:compare-sklearn}

Phép so chỉ có nghĩa khi hai bên tối ưu cùng một hàm. Bộ ước lượng
\texttt{Ridge} của scikit-learn~\cite{pedregosa2011scikit} cực tiểu hóa
$\norm{Xw-y}_{2}^{2} + \alpha \norm{w}_{2}^{2}$ nên quan hệ với
\eqref{eq:objective} là $\alpha = \lambda n$, còn \texttt{SGDRegressor} với
\texttt{penalty="l2"} trùng hàm mục tiêu ngay khi $\alpha = \lambda$. Cách quy
đổi đầy đủ nằm ở phụ lục~\ref{app:conversion}.

Nhóm kiểm phép quy đổi bằng số chứ không bằng suy luận trên giấy: với
$\alpha = \lambda n = \num{1600}$, nghiệm mà \texttt{Ridge} trả về cách nghiệm
đóng \eqref{eq:closed-form} một khoảng $1{,}5 \cdot 10^{-15}$, sai số tương đối
$1{,}9 \cdot 10^{-15}$, và chênh lệch hàm mục tiêu bằng đúng 0 trong số học dấu
phẩy động. Phép kiểm này viết thành một khẳng định trong bộ kiểm thử, vì mọi so
sánh phía sau sẽ vô nghĩa nếu nó thất bại.

\begin{table}[tp]
  \centering
  \caption{Cài đặt mặc định của scikit-learn, chấm bằng hàm mục tiêu
    \eqref{eq:objective} của bài toán. Cột \emph{sai lệch giá} tính theo
    \eqref{eq:rmse-to-price} trên một chiếc máy giá trung vị.}
  \label{tab:sklearn}
  \begin{tabular}{lrrrr}
    \toprule
    Bộ ước lượng & $f(\hat{w}) - \fstar$ & Thời gian (s) & RMSE test & Sai lệch giá \\
    \midrule
    \texttt{Ridge(solver='auto')}     & \num{0}       & \num{0.235} & \num{0.20601} & \num{0}    \\
    \texttt{Ridge(solver='cholesky')} & \num{0}       & \num{0.340} & \num{0.20601} & \num{0}    \\
    \texttt{Ridge(solver='lsqr')}     & \num{2.3e-9}  & \num{0.447} & \num{0.20601} & \num{0}    \\
    \texttt{Ridge(solver='sag')}      & \num{4.9e-9}  & \num{3.524} & \num{0.20601} & \num{0}    \\
    \texttt{SGDRegressor} (mặc định)  & \num{1.16e-3} & \num{0.637} & \num{0.21202} & \num{137}  \\
    \texttt{LinearRegression}         & \num{3.5e-5}  & \num{1.590} & \num{0.20579} & \num{-5}   \\
    \bottomrule
  \end{tabular}
\end{table}

\resultfig{sklearn_iter}
  {So sánh các phương pháp lặp tự cài đặt, theo số vòng lặp.}
  {fig:sklearn-iter}
\resultfig{sklearn_time}
  {So sánh mã tự viết với thư viện, theo thời gian chạy. Các bộ giải trực tiếp
   không có lịch sử vòng lặp nên hiện ra thành một điểm trên mặt phẳng thời gian
   và sai số.}
  {fig:sklearn-time}

Bốn bộ giải \texttt{Ridge} đều đạt $\fstar$ ở mức sai số máy và đều cho cùng một
RMSE, nhưng thời gian trải từ \SI{0.235}{\second} tới \SI{3.52}{\second}
(bảng~\ref{tab:sklearn}). Các bộ giải trực tiếp không có lịch sử vòng lặp nên
hình~\ref{fig:sklearn-iter} chỉ đặt các phương pháp lặp tự cài đặt cạnh nhau, còn
phép so đầy đủ nằm ở hình~\ref{fig:sklearn-time}, nơi mỗi cài đặt thư viện hiện
ra thành một điểm. Bộ giải trực tiếp giải thẳng chính hệ phương trình mà
một bước Newton giải, nên nó và phương pháp Newton tự cài đặt về bản chất làm
cùng một việc; con số \SI{0.133}{\second} của bản tự cài đặt so với
\SI{0.235}{\second} của thư viện không nói lên điều gì ngoài việc bản tự cài đặt
không phải kiểm tra tham số đầu vào và không dựng lại ma trận Gram. Phương pháp
lặp chỉ lấy lại chỗ đứng khi $d$ lớn tới mức không phân rã nổi ma trận
$d \times d$.

Ở tham số mặc định, \texttt{SGDRegressor} dừng cách $\fstar$ một khoảng
$1{,}16 \cdot 10^{-3}$, tức xa hơn 200 lần so với sai số $5{,}8 \cdot 10^{-6}$ mà
cấu hình SGD tự cài đặt đạt được. Hệ quả đo được trên bài toán ứng dụng là
\num{137} đơn vị tiền mỗi máy, con số lớn nhất trong toàn bộ báo cáo và lớn hơn
cả \num{87.7} đơn vị của cấu hình SGD tệ nhất mà nhóm tự cài. Quy tắc chọn độ dài
bước mặc định của bộ ước lượng này không tính theo $L_B$ của bài toán, nên nó rơi
đúng vào lân cận mà chương~\ref{sec:mechanisms} mô tả. Nhận xét chỉ áp cho bộ ước
lượng này ở tham số mặc định, không áp cho thư viện nói chung, và ba dòng
\texttt{Ridge} phía trên là bằng chứng ngược lại.

Dòng cuối bảng đáng chú ý theo hướng ngược lại. \texttt{LinearRegression} không
hiệu chỉnh nên nó không giải bài toán \eqref{eq:objective}, và sai lệch hàm mục
tiêu $3{,}5 \cdot 10^{-5}$ phản ánh đúng điều đó. Nhưng RMSE của nó là
\num{0.20579}, tức tốt hơn nghiệm đóng \num{5} đơn vị tiền, vì quy tắc một sai số
chuẩn ở mục~\ref{sec:choose-lambda} cố ý chọn mức hiệu chỉnh mạnh hơn mức tối ưu
về dự báo. Khoản đổi lấy là số điều kiện nhỏ hơn 100 lần.

\section{Ba câu còn nợ}
\label{sec:compare-open}

Bảng~\ref{tab:summary} trả lời câu hỏi ai về đích trước nhưng chưa giải thích
được câu nào trong ba câu sau, và ba chương tiếp theo nhận lấy từng câu.

Thứ nhất, vì sao ba phương pháp bậc một bám nhau ở đầu rồi tách hẳn ở đuôi, và vì
sao một cấu hình SGD lệch 0,6 đơn vị tiền còn cấu hình kia lệch 87,7. Cả hai đều
quy về ba cơ chế ở chương~\ref{sec:mechanisms}.

Thứ hai, các con số này khớp tới đâu với tốc độ mà lý thuyết dự đoán. Cận co rút
cho gradient descent phụ thuộc $\kappa$, và chương~\ref{sec:theory} đối chiếu hệ
số co rút quan sát được với cận ấy.

Thứ ba, bảng này đo ở đúng một giá trị $\lambda$, tức ở đúng một độ khó. Đổi
$\lambda$ là đổi $\kappa$ theo \eqref{eq:mu-equals-lambda}, và
chương~\ref{sec:lambda} kiểm xem thứ hạng có đổi theo không.
